\documentclass[12pt]{article}
\usepackage{amsmath, amssymb}
\usepackage[margin=1in]{geometry}
\setlength{\parskip}{6pt}
\title{QUBO Formulation: Minimum Multicut on a Tree}
\date{}
\begin{document}
\maketitle

\subsection*{Minimum Multicut on a Tree}

\paragraph{Problem Statement.}
Let $T=(V,E)$ be a tree with edge set $E$ and non-negative edge weights $w_e$ for each $e \in E$. Let $\mathcal{P} = \{(s_1, t_1), \dots, (s_k, t_k)\}$ be a set of $k$ terminal pairs, where $s_i, t_i \in V$. The Minimum Multicut on a Tree problem seeks a subset of edges $C \subseteq E$ such that for every pair $(s_i, t_i)$, the nodes $s_i$ and $t_i$ are disconnected in the graph $T \setminus C$, and the total weight $\sum_{e \in C} w_e$ is minimized.

\paragraph{Decision Variables.}
For each edge $e \in E$, we introduce a binary decision variable $x_e \in \{0,1\}$. The variable $x_e$ indicates whether edge $e$ is selected for the cut: $x_e = 1$ if $e \in C$, and $x_e = 0$ otherwise.

\paragraph{QUBO Derivation.}
\textbf{Step 1: Objective Function.}
The objective is to minimize the total weight of the cut:
\begin{align}
    \min \sum_{e \in E} w_e x_e.
\end{align}

\textbf{Step 2: Constraints.}
For each terminal pair $(s_i, t_i)$, let $P_i \subseteq E$ denote the unique set of edges on the simple path between $s_i$ and $t_i$ in the tree. The separation constraint requires that at least one edge on this path is cut:
\begin{align}
    \sum_{e \in P_i} x_e \ge 1, \quad \forall i \in \{1, \dots, k\}.
\end{align}

\textbf{Step 3: Penalty Terms.}
We convert each inequality constraint to a quadratic penalty term. For a constraint $\sum_{e \in P_i} x_e \ge 1$, the violation occurs when $\sum_{e \in P_i} x_e = 0$. We use the penalty function $\lambda \left(1 - \sum_{e \in P_i} x_e\right)^2$, where $\lambda > 0$ is a penalty weight. Expanding this expression and using the idempotent property $x_e^2 = x_e$ for binary variables:
\begin{align}
    \lambda \left(1 - \sum_{e \in P_i} x_e\right)^2 &= \lambda \left(1 - 2\sum_{e \in P_i} x_e + \left(\sum_{e \in P_i} x_e\right)^2 \right) \\
    &= \lambda \left(1 - 2\sum_{e \in P_i} x_e + \sum_{e \in P_i} x_e^2 + \sum_{e,f \in P_i, e \neq f} x_e x_f \right) \\
    &= \lambda \left(1 - 2\sum_{e \in P_i} x_e + \sum_{e \in P_i} x_e + \sum_{e,f \in P_i, e \neq f} x_e x_f \right) \\
    &= \lambda \left(1 - \sum_{e \in P_i} x_e + \sum_{e,f \in P_i, e \neq f} x_e x_f \right).
\end{align}
The term $\sum_{e,f \in P_i, e \neq f} x_e x_f$ contains each pair $\{e,f\}$ twice (once as $x_e x_f$ and once as $x_f x_e$). Thus, the quadratic contribution for distinct $e,f \in P_i$ is $2\lambda x_e x_f$. The linear contribution for $e \in P_i$ is $-\lambda x_e$, and the constant contribution is $\lambda$.

\textbf{Step 4: QUBO Hamiltonian.}
Combining the objective and all penalty terms, the QUBO Hamiltonian $H(x)$ is:
\begin{align}
    H(x) = \sum_{e \in E} w_e x_e + \sum_{i=1}^k \lambda \left(1 - \sum_{e \in P_i} x_e + \sum_{e,f \in P_i, e \neq f} x_e x_f \right).
\end{align}

\textbf{Step 5: Q Matrix and Offset.}
Collecting coefficients, the QUBO parameters are defined as follows. Let $n_e = |\{i : e \in P_i\}|$ be the number of terminal paths that include edge $e$, and let $n_{e,f} = |\{i : e,f \in P_i\}|$ be the number of terminal paths that include both edges $e$ and $f$. The diagonal entries, off-diagonal entries, and offset are:
\begin{align}
    Q_{e,e} &= w_e - \lambda n_e, \quad \forall e \in E, \\
    Q_{e,f} &= 2\lambda n_{e,f}, \quad \forall e,f \in E, e \neq f, \\
    c &= k \lambda.
\end{align}
The QUBO is then $H(x) = \sum_{e \in E} Q_{e,e} x_e + \sum_{e,f \in E, e < f} Q_{e,f} x_e x_f + c$.

\paragraph{Penalty Weight Justification.}
The penalty weight $\lambda$ must be chosen sufficiently large to ensure that any infeasible solution has a higher energy than the optimal feasible solution. A sufficient condition is $\lambda > \max_{e \in E} w_e$. This ensures that the penalty cost $\lambda$ incurred by violating a separation constraint (which occurs when no edge on a path is cut) exceeds the maximum possible objective saving from omitting any single edge from the cut. Consequently, the optimizer is driven to satisfy all constraints.

\paragraph{Correctness Argument.}
Minimizing $H(x)$ over $x \in \{0,1\}^{|E|}$ recovers the optimal multicut. For any feasible solution, all constraints are satisfied, so the penalty terms vanish and $H(x)$ equals the objective value. For any infeasible solution, at least one constraint is violated, incurring a penalty of at least $\lambda$. Since $\lambda > \max_{e \in E} w_e$, the penalty strictly dominates any potential reduction in the objective function, guaranteeing that the global minimum of $H(x)$ corresponds to a feasible solution with minimum objective value.

\paragraph{Illustrative Example.}
Consider a tree with two edges $E = \{e_0, e_1\}$ and weights $w_0 = 3$, $w_1 = 5$. Let there be a single terminal pair $(s_1, t_1)$ such that the path between them includes both edges, i.e., $P_1 = \{e_0, e_1\}$. Here $k=1$, $n_{e_0} = 1$, $n_{e_1} = 1$, and $n_{e_0, e_1} = 1$. Choosing $\lambda = 10$ (which satisfies $\lambda > \max\{3, 5\}$), the QUBO parameters are:
\begin{align}
    Q_{e_0, e_0} &= 3 - 10(1) = -7, \\
    Q_{e_1, e_1} &= 5 - 10(1) = -5, \\
    Q_{e_0, e_1} &= 2(10)(1) = 20, \\
    c &= 1 \cdot 10 = 10.
\end{align}
The Hamiltonian is:
\begin{align}
    H(x) = -7x_0 - 5x_1 + 20x_0 x_1 + 10.
\end{align}
Evaluating $H(x)$ for all $x \in \{0,1\}^2$:
\begin{itemize}
    \item $x=(0,0)$: $H = 10$ (infeasible, penalty 10).
    \item $x=(1,0)$: $H = -7 + 10 = 3$ (feasible, cost 3).
    \item $x=(0,1)$: $H = -5 + 10 = 5$ (feasible, cost 5).
    \item $x=(1,1)$: $H = -7 - 5 + 20 + 10 = 18$ (feasible, cost 8).
\end{itemize}
The minimum energy is $3$ at $x=(1,0)$, corresponding to cutting edge $e_0$ with total cost $3$.
\end{document}